The goal of this work was to develop
a~library that enables the automatic
optimization of the architecture
of fully connected multilayer perceptrons
using genetic algorithms.
Designing the structure of neural networks
is one of the key stages in the process
of building machine learning models
and it has a~direct impact on their effectiveness,
generalization ability, and training time.
However, the selection of its parameters
- such as the number of layers
or the number of neurons in each layer
- is most often carried out by trial and error,
which requires a~significant investment
of time and expert knowledge.

The theoretical section presents issues
related to artificial neural networks,
with particular emphasis on the MLP architecture.
It also discusses the fundamentals
of evolutionary algorithms,
their mechanisms of operation,
and the application of genetic
algorithms in optimization tasks.
Subsequently, existing solutions used
for the automatic search for neural
network architectures were analyzed
and their limitations were identified.

In the practical section,
the project requirements
were analyzed and based on them,
a~Python library was designed
and described that enables the automatic selection
of a~satisfactory MLP network structure.
The developed solution enables unsupervised
optimization of the architecture based solely
on training and test data,
and the number of input parameters.
It also allows for
the supply of validation data
and the modification of default
solution search parameters,
including the correction term
in the adaptation function.
The library was designed to make it
as easy as possible to:
develop it further
and integrate it into projects.

To evaluate the effectiveness
of the developed solution,
a~series of experiments was
conducted on selected datasets.
The results confirmed the prospect
of effectively using genetic algorithms
to search the vast space of potential
neural network architectures.
The research demonstrated that the library
finds solutions that are at least
as good as those selected manually,
thus confirming the validity
of using genetic algorithms
for designing neural networks.
